% ==========================================
% SECCIÓN 9: ESTRUCTURA METODOLÓGICA
% ==========================================

\section{ESTRUCTURA METODOLÓGICA}

\subsection{Diseño y Nivel de Investigación}

El desarrollo del presente proyecto integrador adopta un \textbf{diseño de investigación aplicada y de desarrollo tecnológico}, caracterizado por un \textbf{nivel descriptivo y propositivo}. La orientación aplicada responde a la necesidad práctica de resolver una problemática concreta de gestión informática en la administración central universitaria. Por su parte, la dimensión tecnológica comprende el diseño e implementación de una arquitectura web modular desacoplada, la normalización de datos relacionales en tercera forma normal, la estructuración de un banco de pruebas de expedientes sintéticos con perfiles ficticios modelados a partir de los formatos institucionales oficiales y el despliegue de una demostración tecnológica interactiva alojada en infraestructura de nube (Vercel y Render), permitiendo validar la solución sin demandar recursos tecnológicos institucionales a la UCE.

\subsection{Métodos de Investigación}

La estrategia metodológica combina métodos teóricos y empíricos fundamentales en la ingeniería de software \cite{pressman_eng}, \cite{sommerville_eng}:

Por una parte, el \textbf{Método Analítico-Sintético} permite descomponer los procesos de gestión de programas doctorales, expedientes estudiantiles, requisitos documentales y el flujo de reservas de citas en tiempo real en sus reglas de negocio constitutivas, para posteriormente sintetizarlas en controladores informáticos, interfaces web desacopladas y servicios Backend centralizados. Por otra parte, el \textbf{Método Inductivo-Deductivo} se emplea para abstraer los requerimientos operacionales observados en las hojas de cálculo locales actuales y deducir los modelos conceptuales y relacionales que rigen el sistema SGD. Finalmente, el \textbf{Método de Ingeniería de Software con Enfoque Ágil Iterativo} articula el ciclo de vida del proyecto en fases concurrentes y sprints quincenales a lo largo de 20 semanas, facilitando la integración continua de componentes, la ejecución de pruebas unitarias y de seguridad en cada entrega, y la validación oportuna de los entregables con los actores de la Unidad.

\subsection{Población de Estudio y Validación Institucional de Requerimientos}

El desarrollo del SGD se fundamenta en una relación de cooperación técnico-académica con la Unidad de Gestión de Doctorados de la Universidad Central del Ecuador, instancia que define las especificaciones técnicas y operativas del sistema. En este contexto institucional, para garantizar la máxima rigurosidad metodológica, la independencia empírica y la ética de la investigación, resulta imperativo prevenir cualquier conflicto de interés metodológico (situación de ``juez y parte'') que pudiera derivarse del hecho de que la dirección de la Unidad coincida con la tutoría académica del presente proyecto de titulación. Por consiguiente, mientras el Director de la Unidad ejerce estrictamente su rol de dirección institucional y tutoría metodológica sin intervenir como sujeto experimental, la \textbf{población y muestra de evaluación operativa y pruebas de aceptación (\textit{User Acceptance Testing} - UAT)} queda conformada de forma autónoma e independiente por los $N = 2$ analistas técnicos y operativos de la Unidad: la \textit{Analista de Procesos Universitarios} y el \textit{Analista Técnico de Sistemas / Gestión Académica}. Ellos constituyen los usuarios operativos directos competentes para emitir el dictamen formal y técnico de aceptación sobre la funcionalidad del prototipo.

Por su parte, los docentes titulares becarios de doctorado corresponden al universo objetivo del dominio de negocio sobre el cual se modelan las reglas contractuales, plazos y devengamientos, pero \textbf{no actúan como sujetos de prueba ni evaluadores del prototipo}, garantizando que el proceso de validación no dependa de la disponibilidad de los estudiantes ni introduzca sesgos operacionales. Para salvaguardar la privacidad institucional conforme a la LOPDP y asegurar que ninguna persona real se encuentre expuesta en el entorno público de pruebas, la base de datos de la demo se puebla mediante un \textbf{banco de datos sintéticos con perfiles e identidades ficticias de ejemplo}. Estos registros fueron generados tomando como base estricta las estructuras, casuísticas y tipos de trámite de los archivos de Excel oficiales de la Unidad, simulando fielmente los escenarios administrativos reales (prórrogas, adendas y suspensiones temporales) pero sin contener información de personas reales.

Bajo este marco, la evaluación final de la plataforma prescinde de encuestas de usabilidad y se canaliza al 100\% a través de la \textbf{Matriz de Validación y Conformidad de Requerimientos Funcionales y No Funcionales} \cite{ieee_requirements}, suscrita de manera independiente por los dos analistas técnicos y operativos de la Unidad ($N = 2$), estableciéndose como criterio óptimo de éxito y aprobación institucional alcanzar una tasa de completitud funcional igual o superior al 90\% ($\ge 90\%$).

\subsection{Técnicas e Instrumentos de Recolección de Datos}

Para el levantamiento empírico de requerimientos y la evaluación de la solución, se emplean tres técnicas e instrumentos principales:

En primer lugar, se aplica la \textbf{Entrevista Semiestructurada} orientada al levantamiento inicial de requerimientos y dinámicas institucionales, en la cual participa el personal directivo y técnico de la Unidad, utilizando como instrumento una \textit{Guía de Entrevista Cualitativa} enfocada en mapear los flujos de trabajo, las excepciones administrativas y las fórmulas de cómputo de adendas. En segundo lugar, se recurre al \textbf{Análisis Documental}, empleando una \textit{Matriz de Extracción de Datos} para examinar los formatos de las hojas de cálculo de Excel históricas y las cláusulas del Reglamento de Régimen Académico UCE. En tercer lugar, para la evaluación empírica final del sistema, se aplica una \textbf{Matriz de Validación y Conformidad de Requerimientos Funcionales y No Funcionales} \cite{ieee_requirements}, instrumento técnico suscrito de forma independiente por los analistas técnicos y operativos ($N = 2$) que evalúa el grado de completitud operativa, la exactitud en el cómputo de prórrogas y adendas, la eficacia del módulo de solicitudes y citas, y la adecuación de los manuales interactivos entregados.

\subsection{Procesamiento de Datos y Verificación de Hipótesis}

Los datos recopilados durante las fases de diagnóstico, pruebas de software y evaluación serán procesados sistemáticamente. A partir del análisis de las hojas de cálculo históricas oficiales, se estructurará el banco de datos sintético con perfiles ficticios de prueba, el cual será ingerido en la base de datos PostgreSQL mediante Prisma ORM para garantizar integridad referencial y consistencia de tipos. Posteriormente, la lógica de cálculo contractual, el generador de certificados digitales y el motor de alertas se programarán en Express.js respaldados por pruebas unitarias automatizadas en Jest y pruebas de integración con Supertest.

La verificación de la hipótesis se realizará contrastando el cómputo automatizado frente al 100\% de los casos de prueba teóricos de devengamiento y prórrogas. Por otra parte, al finalizar el desarrollo, la Unidad de Gestión de Doctorados evaluará la plataforma mediante sesiones de prueba de aceptación (\textit{User Acceptance Testing} - UAT), en las cuales los analistas técnicos y operativos ($N = 2$), de forma metodológicamente independiente del docente tutor, verificarán exhaustivamente los módulos implementados frente a la especificación de requisitos del sistema (SRS), midiendo la tasa de completitud funcional (con meta $\ge 90\%$) como criterio determinante de conformidad y éxito para la entrega del proyecto de titulación.

\subsection{Fases del Desarrollo del Proyecto (Planificación de 20 Semanas)}

El ciclo de vida del proyecto se articula en cinco fases metodológicas consecutivas a lo largo de un horizonte de 20 semanas (5 meses):

\begin{enumerate}[label=\textbf{Fase \arabic*:},leftmargin=*]
    \item \textbf{Análisis de Esquemas, Generación de Datos Sintéticos y Requisitos (Mes 1, Semanas 1 a 4):} Levantamiento de requerimientos con el personal de la Unidad, análisis de las hojas de cálculo Excel oficiales, estructuración algorítmica del banco de expedientes sintéticos con perfiles ficticios de prueba y formulación del documento de requisitos (SRS).
    \item \textbf{Diseño de Arquitectura, Modelado Prisma y Prototipado (Mes 2, Semanas 5 a 8):} Diseño del esquema relacional en PostgreSQL mediante Prisma ORM, especificación de la matriz de permisos RBAC, definición de contratos de interfaz desacoplados y maquetación visual con TailwindCSS.
    \item \textbf{Desarrollo Backend REST, Pruebas Unitarias y Frontend Admin (Meses 3 y 4, Semanas 9 a 16):} Construcción de servicios API REST en Express.js con tokens JWT firmados, desarrollo del Panel Administrativo en Next.js 16 (Bun/Zustand), y ejecución de pruebas unitarias automatizadas (Jest) y de seguridad OWASP.
    \item \textbf{Desarrollo Portal Estudiantil, Sistema de Reservas de Citas y Alertas (Meses 4 y 5, Semanas 13 a 18):} Construcción del Portal Estudiantil en Astro y React, implementación del visor de estado del expediente (documentos entregados y faltantes), desarrollo del sistema de reservas de citas en tiempo real con el personal de la Unidad, generador de certificados digitales y motor preventivo de alertas de vencimiento.
    \item \textbf{Despliegue Cloud, Documentación Fumadocs y Validación Institucional (Mes 5, Semanas 17 a 20):} Despliegue de la demo funcional en la nube (Vercel y Render), compilación de manuales interactivos en Fumadocs (Manual de Usuario y Guía Técnica de Arquitectura), aplicación de la Matriz de Validación de Requerimientos de forma independiente con los analistas de la Unidad ($N = 2$, criterio de éxito $\ge 90$\%) y elaboración del informe final de titulación.
\end{enumerate}
